\documentclass[10pt,a4paper]{amsart}
\usepackage[margin=19mm]{geometry}
\usepackage{amsmath,amssymb,mathtools}
\usepackage[scr=rsfs,cal=euler]{mathalfa}
\usepackage{xfrac}
\usepackage[unicode,hidelinks,bookmarks=false]{hyperref}
\newcommand{\ie}{i.e.\ }
\newcommand{\cf}{cf.\ }
\newcommand{\resp}{respectively\ }
\newcommand{\Subsection}{\subsection}
\providecommand{\nobreakdash}{\nobreak\hbox{-}}
\setlength{\emergencystretch}{2em}
\multlinegap0pt
\usepackage{mathtools}
\usepackage[shortlabels]{enumitem}
\usepackage{amssymb}
\usepackage{bm}
\usepackage{tikz}
\usepackage{float}
\newtheorem{proposition}{Proposition}
\title{Nonlinear Dynamics In Optimization Landscape of Shallow Neural Networks with Tunable Leaky ReLU}
\author{Jingzhou Liu}
\date{Source: arXiv:2510.25060v2 (CC BY 4.0)}
\begin{document}
\maketitle

\noindent\emph{Source and licence.} Nonlinear Dynamics In Optimization Landscape of Shallow Neural Networks with Tunable Leaky ReLU. Version arXiv:2510.25060v2 (CC BY 4.0), distributed by arXiv under the Creative Commons Attribution 4.0 International licence, http://creativecommons.org/licenses/by/4.0/, as recorded in the arXiv licence metadata for this exact version and shown on the abstract page https://arxiv.org/abs/2510.25060v2. Full licence notice: \texttt{licenses/arxiv-2510.25060-CC-BY-4.0.txt}.\par\smallskip

\noindent\textit{Editorial scope.} Counted unit: the article's proposition frobenius (source0.tex lines 1010--1020). The article's complete original proof of it is delivered with it (source0.tex lines 1021--1069, one \texttt{proof} environment of 49 lines). No mathematics is written, repaired or re-derived: the statement and the proof are the article's own lines, in the article's own notation, and no retained line is substituted or re-flowed.  Retained uncounted as the source-local context the proof needs: lines 984--991.  Nothing was omitted from any retained span, so the package carries no omission record.  No retained line contains a citation, so the wrapper carries no bibliography and no bibliography entry.  No retained line contains a figure environment, an image-inclusion command or drawing code, so no image is delivered and the package declares no asset dependency.  Editorial formatting only, in addition: an AMS \texttt{amsart} preamble that restates the article's own declarations and macros used by the retained lines, namely the environments \texttt{proposition} on separate counters, together with the standard packages those lines need.  The retained environments print in this document as Proposition 1 in order of appearance, whereas the article prints the same items with the numbering of its own full document.  The complete native source of the article as distributed in the arXiv e-print for this exact version is bundled unchanged as \texttt{sources/188003/source0.tex}.
\begin{equation}\label{eq:frobenius}
\chi_\eta(C_i)
= 
\bigl[
\Delta(x)\,
\prod_{j} P_j(x)^{\,i_j}
\bigr]_{(l_1,\dots,l_r)},
\end{equation}

\begin{proposition}\label{prop: frobenius} Let $k = \sum_j j\,i_j$, then we have
  \begin{enumerate}
      \item[(a)] If $\eta = (k),$ then $\chi_{(k)}(C_i) = 1.$
      \item[(b)] If $\eta = (k-1,1),$ then $\chi_{(k-1,1)}(C_i) = i_1 - 1.$ 
      \item[(c)] If $\eta= (k-2,1,1)$, then $\chi_{(k-2,1,1)}(C_i)=\frac{1}{2}(i_1-1)(i_1-2) - i_2.$
      \item[(d)] If $\eta = (k-2,2)$,
$\chi_{(k-2,2)}(C_i)
=
\frac{1}{2}(i_1-1)(i_1-2) + i_2 - 1.$
  \end{enumerate}  
\end{proposition}

\begin{proof} By applying Frobenius Formula \eqref{eq:frobenius}, we have
\begin{enumerate}
\item[(a)] For $\eta = (k),$ we have $r=1$, $l_1 = k$ and $\Delta(x)=1$, $P_j(x)=x_1^j$.
Thus, $\chi_{(k)}(C_i) = 1$ for all $i$.
\item[(b)] For $\eta = (k-1,1)$. Then $r=2$ and $l=(k,1).$ Thus,
\[
\chi_{(k-1,1)}(C_i)
=
\Bigl[
(x_1-x_2)\,
\prod_j (x_1^j + x_2^j)^{i_j}
\Bigr]_{x_1^{k}x_2^{1}}.
\]
Let $\Phi(x_1,x_2):=\Bigr[\prod_j (x_1^j + x_2^j)^{i_j}
\Bigr]_{x_1^{k}x_2^{1}},$ then $
\chi_{(k-1,1)}(C_i)=\Bigr[x_1\Phi\Bigr]_{x_1^{k}x_2^{1}}-\Bigr[x_2\Phi\Bigr]_{x_1^{k}x_2^{1}}$ which is equivalent to 
\[
\Bigr[\Phi\Bigr]_{x_1^{k-1}x_2^{1}}-\Bigr[\Phi\Bigr]_{x_1^{k}x_2^{0}}.
\]
For $\Bigr[\Phi\Bigr]_{x_1^{k}x_2^{0}},$ one has for each $j,$ elements of $(x_1^j+x_2^j)^{i_j}$ are of the form 
\[
(x_1^j)^{a_j}(x_2^j)^{b_j},\;\;a_j+b_j=i_j.
\]
Therefore, 
\[
\Bigr[\Phi\Bigr]_{x_1^{k}x_2^{0}}=\Bigr[\prod_{j}\binom{i_j}{a_j}x_1^{\sum_j ja_j}x_2^{\sum_j jb_j}\Bigr]_{x_1^kx_2^0},
\]
which implies
\[
\sum_j ja_j = k,\quad
\sum_j jb_j=0.
\]
i.e. $b_j=0$ and $a_j=i_j,$ so $\Bigr[\Phi\Bigr]_{x_1^{k}x_2^{0}}=\binom{i_j}{i_j}=1.$ Similarly, $\Bigr[\Phi\Bigr]_{x_1^{k-1}x_2^{1}}=i_1.$
Therefore,
\[
\chi_{(k-1,1)}(C_i) = i_1 - 1.
\]
\item[(c)] For $\eta= (k-2,1,1)$, we have $r=3$ with $l = (k,2,1)$, and
\[
\chi_{(k-2,1,1)}(C_i)
=
\Bigl[
(x_1-x_2)(x_1-x_3)(x_2-x_3)
\prod_j (x_1^j + x_2^j + x_3^j)^{i_j}
\Bigr]_{x_1^{k}x_2^{2}x_3^{1}}.
\]
\item[(d)] For $\eta = (k-2,2),$ similarly, we have $r=2$ with $l = (k-1,2)$.
\end{enumerate}
\end{proof}

\end{document}
